% Part: incompleteness
% Chapter: incompleteness-provability
% Section: lob-thm

\documentclass[../../../include/open-logic-section]{subfiles}

\begin{document}

% SEGMENT OLP-0320-S01

\olfileid{inc}{inp}{lob}

\olsection{லோபின் தேற்றம்}

ஒரு கோட்பாடு~$\Th{T}$ க்கான கோடல் வாக்கியம், $\lnot
\OProv[\Th{T}](y)$ இன் நிலைப்புள்ளியாகும்; அதாவது,
\[
\Th{T} \Proves \lnot \OProv[\Th{T}](\gn{!G}) \liff !G.
\]
என்பதை நிறைவேற்றும் !!a{sentence}~$!G$ ஆகும். இது
!!{derivable} அல்ல. ஏனெனில் $\Th{T} \Proves !G$ என்றால்,
(a) !!{derivability} நிபந்தனை (1) இன் படி
$\Th{T} \Proves \OProv[\Th{T}](\gn{!G})$; மேலும் (b)
$\Th{T} \Proves !G$ மற்றும்
$\Th{T} \Proves \lnot \OProv[\Th{T}](\gn{!G}) \liff !G$
ஆகியவற்றிலிருந்து
$\Th{T} \Proves \lnot \OProv[\Th{T}](\gn{!G})$ கிடைக்கும்; எனவே
$\Th{T}$ முரணாகிவிடும். இப்போது
$\OProv[\Th{T}](y)$ இன் நிலைப்புள்ளி, அதாவது
\[
\Th{T} \Proves \OProv[\Th{T}](\gn{!H}) \liff !H.
\]
என்பதை நிறைவேற்றும் !!a{sentence}~$!H$ பற்றிக் கேட்பது இயல்பு.
இது !!{derivable} என்றால், நிபந்தனை (1) மூலம்
$\Th{T} \Proves \OProv[\Th{T}](\gn{!H})$ கிடைக்கும்; ஆனால் மேலுள்ள
சமன்பாட்டிலிருந்து modus ponens பயன்படுத்தினாலும் அதே முடிவு கிடைக்கும்.
எனவே, குறைந்தபட்சம் கோடல் வாக்கியத்தின் வழக்கில் பயன்படுத்திய அதே
வாதத்தால் $\Th{T}$ முரணானது என்று நாம் பெற முடியாது. இதனால்
$\Th{T}$ \emph{உண்மையில்} $!H$ ஐ !!{derive} செய்கிறது என்று தெரியவில்லை.

இந்தக் கேள்வியைச் சிறிது பொதுமைப்படுத்தினால் முன்னேறலாம். நிலைப்புள்ளி
சமன்பாட்டின் இடமிருந்து வலமாகச் செல்லும் திசையான
$\OProv[\Th{T}](\gn{!H}) \lif !H$ என்பது, பொதுத் திட்டத்தின் ஓர்
நிகழ்வாகும்; இதை \emph{பிரதிபலிப்புக் கொள்கை} என்று அழைக்கிறோம்:
$\OProv[\Th{T}](\gn{!A}) \lif !A$. $\Th{T}$ தன்னால் என்ன
!!{derive} செய்ய முடியும் என்பதைப் பற்றிப் “பிரதிபலிக்கிறது” என்ற
பொருளில் இது அப்பெயரைப் பெறுகிறது; அடிப்படையில், “$\Th{T}$,
$!A$ ஐ !!{derive} செய்ய முடிந்தால், $!A$ மெய்யாகும்” என்று அது
$!A$ எதற்கும்
கூறுகிறது. இது மெய்மையான கோட்பாடுகளுக்கு மட்டுமே பொருந்தும்; ஆகவே
பொதுவாக ஒரு கோட்பாடு அதன் எல்லா நிகழ்வுகளையும் !!{derive} செய்யாது
என்பது இயல்பான எதிர்பார்ப்பு. அப்படியானால், போதிய வலிமை கொண்டதும்
!!{derivability} நிபந்தனைகளை நிறைவேற்றுவதுமான ஒரு கோட்பாடு எந்த
நிகழ்வுகளை !!{derive} செய்ய முடியும்? $!A$ தானே !!{derivable} ஆக
இருக்கும் நிகழ்வுகள் அனைத்தையும் அது நிச்சயமாகச் செய்யும். அவ்வளவுதான்,
அடுத்த முடிவு காட்டுவது இதுவே.

\begin{thm}\ollabel{thm:lob}
$\Th{Q}$ ஐ நீட்டிக்கும் !!a{axiomatizable} கோட்பாடு~$\Th{T}$ ஐ
எடுத்துக்கொள்க. மேலும் $\OProv[\Th{T}](y)$ என்பது
\olref[2in]{sec} இல் உள்ள P1--P3 நிபந்தனைகளை நிறைவேற்றும் வாய்பாடு
என்க. $\Th{T}$,
$\OProv[\Th{T}](\gn{!A}) \lif !A$ ஐ !!{derive} செய்கிறது எனில்,
உண்மையில் $\Th{T}$, $!A$ ஐ !!{derive} செய்கிறது.
\end{thm}

வேறு விதமாகச் சொன்னால், $\Th{T} \Proves/ !A$ எனில்,
$\Th{T} \Proves/
\OProv[\Th{T}](\gn{!A}) \lif !A$. இந்த முடிவு லோபின் தேற்றம்
என்று அழைக்கப்படுகிறது.

\begin{explain}
லோபின் தேற்றத்தின் நிறுவலுக்கான உள்ளுணர்வு, சாந்தா கிளாஸ்
இருக்கிறார் என்பதை நிரூபிக்கும் ஒரு நுட்பமான வாதமாகும். அந்த முடிவு
உங்களுக்குப் பிடிக்கவில்லை என்றால், நீங்கள் விரும்பும் வேறு எந்த
முடிவையும் இங்கு மாற்றிக் கொள்ளலாம். வாதம் இதோ:
\begin{enumerate}
\item $X$ என்பது “$X$ மெய்யானால் சாந்தா கிளாஸ் இருக்கிறார்” என்ற
  வாக்கியம்.
\item $X$ மெய்யானது என்று கருதுக.
\item அப்போது அந்த வாக்கியம் கூறுவது பொருந்தும்; அதாவது, $X$ மெய்யானால் சாந்தா
  கிளாஸ் இருக்கிறார்.
\item $X$ மெய்யானது என்று நாம் கருதுவதால், (2), (3) இலிருந்து
  modus ponens மூலம் சாந்தா கிளாஸ் இருக்கிறார் என்று முடிவு செய்யலாம்.
\item “$X$ மெய்யானது” என்ற அனுமானத்திலிருந்து (4), “சாந்தா கிளாஸ்
  இருக்கிறார்” என்பதை நாம் பெற்றுவிட்டோம். நிபந்தனை நிறுவலின் மூலம்,
  “$X$ மெய்யானால் சாந்தா கிளாஸ் இருக்கிறார்” என்று காட்டியுள்ளோம்.
\item ஆனால் இது $X$ என்ற வாக்கியமே. எனவே $X$ மெய்யானது என்று
  காட்டியுள்ளோம்.
\item அப்போது (2)--(4) என்ற வாதத்தின் மூலம் சாந்தா கிளாஸ் இருக்கிறார்.
\end{enumerate}
“மெய்யானது” என்பதற்குப் பதிலாக “!!{derivable}” என்றும், “சாந்தா
கிளாஸ் இருக்கிறார்” என்பதற்குப் பதிலாக~$!A$ என்றும் வைத்து இந்த
எண்ணத்தை முறைப்படுத்தினால் லோபின் தேற்றத்தின் நிறுவல் கிடைக்கும்.
இதில் உள்ள தந்திரம், !!{formula}~$\OProv[\Th{T}](y) \lif !A$ க்கு
நிலைப்புள்ளித் துணைத்தேற்றத்தைப் பயன்படுத்துவதாகும். அதன் நிலைப்புள்ளி
முந்தைய வரைவிலுள்ள $X$ வாக்கியத்துக்கு ஒத்ததாக இருக்கும்.
\end{explain}

\begin{proof}[Proof of \olref{thm:lob}]
$!A$ என்பது !!a{sentence} என்றும், $\Th{T}$ ஆனது
$\OProv[\Th{T}](\gn{!A}) \lif !A$ ஐ !!{derive}s செய்கிறது என்றும்
கருதுக. $!B(y)$ என்பது !!{formula}~$\OProv[\Th{T}](y) \lif !A$
ஆக இருக்கட்டும். நிலைப்புள்ளித் துணைத்தேற்றத்தைப் பயன்படுத்தி,
$\Th{T}$ ஆனது $!D \liff !B(\gn{!D})$ ஐ !!{derive}s செய்யுமாறு
அமையும் !!a{sentence}~$!D$ ஒன்றைக் கண்டுபிடிப்போம்.
அப்போது பின்வரும் ஒவ்வொன்றும் $\Th{T}$ இல் !!{derivable} ஆகும்:
\begin{align}
  & !D \liff (\OProv[\Th{T}](\gn{!D}) \lif !A) \ollabel{L-1}\\
  & \qquad \text{$!D$ என்பது~$!B(y)$ இன் நிலைப்புள்ளி} \notag \\
  & !D \lif (\OProv[\Th{T}](\gn{!D}) \lif !A) \ollabel{L-2}\\
  & \qquad \text{\olref{L-1} இலிருந்து} \notag\\
  & \OProv[\Th{T}](\gn{!D \lif (\OProv[\Th{T}](\gn{!D}) \lif !A)}) \ollabel{L-3}\\
  & \qquad \text{\olref{L-2} இலிருந்து P1 நிபந்தனையின் மூலம்} \notag \\
  & \OProv[\Th{T}](\gn{!D}) \lif \OProv[\Th{T}](\gn{\OProv[\Th{T}](\gn{!D}) \lif !A})
  \ollabel{L-4}\\
  &\qquad \text{\olref{L-3} இலிருந்து P2 நிபந்தனையைப் பயன்படுத்தி} \notag \\
  & \OProv[\Th{T}](\gn{!D}) \lif (\OProv[\Th{T}](\gn{\OProv[\Th{T}](\gn{!D})}) \lif \OProv[\Th{T}](\gn{!A})) \ollabel{L-5}\\
  &\qquad \text{\olref{L-4} இலிருந்து மீண்டும் P2 நிபந்தனையைப் பயன்படுத்தி} \notag\\
& \OProv[\Th{T}](\gn{!D}) \lif \OProv[\Th{T}](\gn{\OProv[\Th{T}](\gn{!D})}) \ollabel{L-6}\\
  & \qquad\text{!!{derivability} நிபந்தனை P3 மூலம்} \notag\\
  & \OProv[\Th{T}](\gn{!D}) \lif \OProv[\Th{T}](\gn{!A}) \ollabel{L-7} \\
  &\qquad\text{\olref{L-5}, \olref{L-6} ஆகியவற்றிலிருந்து}\notag\\
  & \OProv[\Th{T}](\gn{!A}) \lif !A \ollabel{L-8}\\
  &\qquad\text{தேற்றத்தின் அனுமானத்தின்படி} \notag\\
  & \OProv[\Th{T}](\gn{!D}) \lif !A \ollabel{L-9}\\
  &\qquad\text{\olref{L-7}, \olref{L-8} ஆகியவற்றிலிருந்து}\notag\\
  & (\OProv[\Th{T}](\gn{!D}) \lif !A) \lif !D \ollabel{L-10}\\
  & \qquad \text{\olref{L-1} இலிருந்து} \notag \\
  & !D \ollabel{L-11}\\
  & \qquad\text{\olref{L-9}, \olref{L-10} ஆகியவற்றிலிருந்து}\notag \\
  & \OProv[\Th{T}](\gn{!D}) \ollabel{L-12}\\
  & \qquad\text{\olref{L-11} இலிருந்து P1 நிபந்தனையின் மூலம்}\notag \\
  & !A \qquad\qquad\text{\olref{L-8}, \olref{L-12} ஆகியவற்றிலிருந்து}\notag
\end{align}
\end{proof}

லோபின் தேற்றம் கையில் இருந்தால், (P1)--(P3) நிபந்தனைகளை நிறைவேற்றும்
!!a{derivability} பயனிலை கொண்ட கோட்பாடுகளுக்கான இரண்டாம்
முழுமையின்மைத் தேற்றத்துக்கு ஒரு குறுகிய நிறுவல் கிடைக்கிறது:
$\Th{T} \Proves \OProv[\Th{T}](\gn{\lfalse}) \lif \lfalse$ எனில்,
$\Th{T} \Proves \lfalse$. $\Th{T}$ முரணற்றது என்றால்,
$\Th{T} \Proves/ \lfalse$. ஆகவே
$\Th{T} \Proves/ \OProv[\Th{T}](\gn{\lfalse}) \lif \lfalse$,
அதாவது $\Th{T} \Proves/ \OCon[\Th{T}]$. மேலும்,
$\OProv[\Th{T}](x)$ இன் நிலைப்புள்ளியான $!H$ !!{derivable} என்பதை
நிறுவவும் இதைப் பயன்படுத்தலாம். ஏனெனில்
\begin{align*}
  \Th{T} & \Proves \OProv[\Th{T}](\gn{!H}) \liff !H\\
  \intertext{குறிப்பாக}
    \Th{T} & \Proves \OProv[\Th{T}](\gn{!H}) \lif !H
\end{align*}
என்பதால், லோபின் தேற்றத்தின் மூலம் $\Th{T} \Proves !H$.

\begin{prob}
$\Th{T}$ ஒரு கணிக்கத்தக்கவாறு அடிகோளாக்கப்பட்ட கோட்பாடு என்றும்,
$\OProv[\Th{T}]$ என்பது $\Th{T}$ க்கான !!a{derivability} பயனிலை
என்றும் கொள்க. பின்வரும் நான்கு கூற்றுகளைக் கருதுக:
\begin{enumerate}
\item $T \Proves !A$ எனில், $T \Proves
  \OProv[\Th{T}](\gn{!A})$.
\item $T \Proves !A \lif \OProv[\Th{T}](\gn{!A})$.
\item $T \Proves \OProv[\Th{T}](\gn{!A})$ எனில்,
  $T \Proves !A$.
\item $T \Proves \OProv[\Th{T}](\gn{!A}) \lif !A$.
\end{enumerate}
ஒவ்வொரு கூற்றும் எந்த நிபந்தனைகளின் கீழ் உண்மையாகும்?
\end{prob}

\end{document}
